\scanpage{620}
\mspage{8}
Un couple $Q$-orthogonal \inserted{$\Phi\longleftrightarrow\Psi$} est dit \deleted{\unsure{$Q$|s}} \uline{spécial}
s'il satisfait les conditions suivantes

\[
\left\{\begin{aligned}
\text{a)}\quad &M\ \text{est \unsure{stable|stble} par \inserted{petites} }\unsure{\underline{\lim}|\varinjlim},\ \text{les objets de }M\text{ sont}\\
&\text{accessibles}\quad [\text{On dit alors que }M\text{ est ``\unsure{spécial|spéciale}''}]\\
\text{b)}\quad &\deleted{\Phi\ \text{\unsure{est}}}\ \text{Le couple est clos : }\Phi\longleftrightarrow\Psi.\\
\text{c)}\quad &\Phi\ \text{est \uline{spéciale}.}
\end{aligned}\right.\tag*{(1.10)}
\]

Cela implique donc, par 1.1, que $\Phi \mathrel{\longleftarrow\!\mid\!\longrightarrow} \Psi$ est
un couple de Quillen. Si $M$ satisfait a),
les couples $Q$-orthogonaux spéciaux dans $M$,
ou encore les couples de Quillen spéciaux,
correspondent donc biunivoquement aux
parties $\Phi$ de $\mathrm{Fl}\,M$ qui sont \uline{spéciales} et
closes à gauche (ou ce qui revient au $=$, $Q$-saturées
à gauche), en \deleted{\unsure{fai}} associant à une telle
$\Phi$ le couple $\Phi \mathrel{\longleftarrow\!\mid\!\longrightarrow} \Phi_{\unsure{*|\alpha}}$.

Ceci posé :

\uline{Corollaire} 1.8 : Supposons que $M,M'$ satisfassent
\margin{\tikz[x=1em,y=1em]{\draw (0,0) .. controls (-.08,-1.3) and (.12,-2.6) .. (.22,-3.6); \draw (.065,.06) .. controls (.035,-1.4) and (.25,-2.5) .. (.12,-3.45);}}
la condition a) de (1.10), et soit $(\Phi,\Psi)$
un couple de Quillen spécial dans $M$. Alors
$\big(\Phi'=u_!(\Phi)^{\sim},\ \Psi'=(u^*)^{-1}(\deleted{\unsure{M}}\Psi)\big)$ est un
couple de Quillen spécial dans $M'$.

En effet, par le cor. 1.6 on a

\[
\Phi'_*=\big(u_!(\Phi)\big)_{\unsure{*|\alpha}}=(u^{*})^{-1}(\Psi)=\Psi',
\]

donc $\Phi'\longleftrightarrow\Psi'$. Il reste à prouver que
$\Phi'$ est \unsure{spécial|spéciale}. Soit $\Phi_0$ une partie \unsure{petite} de
$\Phi$ telle que $\Phi\subset\Phi_0^{\sim}$, posons $\Phi_0'=u_!(\Phi_0)\inserted{\subset}\Phi'$
% archive sheet number: 619
